% English translation of chapter1.tex, lines 568--629.
% Source: Methods of Algebra, Volume 2, commit
% 9a5803ff77dd3257484cb177f851a73770a59dd3 (CC BY 4.0).
% stable-unit-id: o014.aljabr2.chapter1.linear-categories-over-commutative-rings

% segment-id: o014.li2.chapter1.linear-categories-over-commutative-rings.h001
\section{Generalization: Linear Categories over a Commutative Ring}\label{sec:linear-cat}
\hypertarget{o014.aljabr2.chapter1.linear-categories-over-commutative-rings}{ }

% segment-id: o014.li2.chapter1.linear-categories-over-commutative-rings.p001
The $\Hom$ sets of an additive category have the structure of abelian groups.
In many situations they also carry scalar multiplication by some commutative
ring $\Bbbk$. For example, every $\Hom$ set in $\Bbbk\dcate{Mod}$ is naturally
a $\Bbbk$-module, and composition of morphisms is $\Bbbk$-bilinear.

% segment-id: o014.li2.chapter1.linear-categories-over-commutative-rings.d001
\begin{definition}\label{def:cat-k-linear}
	\index{category!$\Bbbk\dcate{Mod}$-}
	\index{linear!$\Bbbk$- ($\Bbbk$-linear)}
	Let $\Bbbk$ be a commutative ring. Suppose that every $\Hom$ set in a
	category $\mathcal{A}$ is equipped with a $\Bbbk$-module structure such that
	composition of morphisms
	$\Hom(Y, Z) \times \Hom(X, Y) \to \Hom(X, Z)$ is a $\Bbbk$-bilinear map
	for all objects $X,Y,Z$; in other words, the following diagram commutes:
	% segment-id: o014.li2.chapter1.linear-categories-over-commutative-rings.g001
	\[\begin{tikzcd}
		\Hom(Y, Z) \times \Hom(X, Y) \arrow[d] \arrow[r, "\text{composition}"] & \Hom(X, Z) \\
		\Hom(Y, Z) \dotimes{\Bbbk} \Hom(X, Y) \arrow[ru, bend right, start anchor=east, "{\exists! \; \text{$\Bbbk$-module homomorphism}}"']
	\end{tikzcd}  \]
	With these data, $\mathcal{A}$ is called a \emph{$\Bbbk\dcate{Mod}$-category}.

	Let $F: \mathcal{A} \to \mathcal{A}'$ be a functor between
	$\Bbbk\dcate{Mod}$-categories. If the map
	$\Hom(X,Y) \to \Hom(FX,FY)$ is a $\Bbbk$-module homomorphism for every
	$X,Y \in \Obj(\mathcal{A})$, then $F$ is called a
	\emph{$\Bbbk$-linear functor}.
\end{definition}

% segment-id: o014.li2.chapter1.linear-categories-over-commutative-rings.p002
Observe that $\Bbbk\dcate{Mod}$ is a monoidal category under the tensor product
$\otimes := \dotimes{\Bbbk}$. Thus the term ``$\Bbbk\dcate{Mod}$-category''
agrees with the terminology for enriched categories in \cite[\S 3.4]{Li1}.
The $\cate{Ab}$-categories introduced earlier are precisely the
$\Z\dcate{Mod}$-categories. Conversely, forgetting scalar multiplication on a
$\Bbbk\dcate{Mod}$-category leaves an $\cate{Ab}$-category.

% segment-id: o014.li2.chapter1.linear-categories-over-commutative-rings.e001
\begin{example}
	A $\Bbbk\dcate{Mod}$-category with a single object is nothing but a
	$\Bbbk$-algebra. The correspondence is given by
	$\mathcal{A} \mapsto \End_{\mathcal{A}}(\star)$, where
	$\Obj(\mathcal{A}) = \{\star\}$.
\end{example}

% segment-id: o014.li2.chapter1.linear-categories-over-commutative-rings.p003
The argument of Proposition~\ref{prop:automatic-Ab-morphism} carries over
verbatim and gives the following result.

% segment-id: o014.li2.chapter1.linear-categories-over-commutative-rings.pr001
\begin{proposition}\label{prop:automatic-k-morphism}
	Let $\mathcal{A}$ and $\mathcal{A}'$ be $\Bbbk\dcate{Mod}$-categories.
	% segment-id: o014.li2.chapter1.linear-categories-over-commutative-rings.l001
	\begin{enumerate}
		% segment-id: o014.li2.chapter1.linear-categories-over-commutative-rings.i001
		\item Given a pair of adjoint functors
		% The originally inline source adjunction diagram was reflowed as a
		% centered display so that it neither crowds nor overruns the reader width.
		% The long chain of equivalences in the proof was likewise reflowed.
		% segment-id: o014.li2.chapter1.linear-categories-over-commutative-rings.g002
		\[
		\begin{tikzcd}
			F: \mathcal{A} \arrow[shift left, r] & \mathcal{A}' \arrow[shift left, l] : G
		\end{tikzcd},
		\]
		with both $F$ and $G$ $\Bbbk$-linear, the adjunction isomorphism
		$\Hom_{\mathcal{A}'}(F(\cdot),\cdot) \rightiso
		\Hom_{\mathcal{A}}(\cdot,G(\cdot))$ is $\Bbbk$-linear.
		% segment-id: o014.li2.chapter1.linear-categories-over-commutative-rings.i002
		\item For every $\alpha: I \to \mathcal{A}$, the bijection
		% segment-id: o014.li2.chapter1.linear-categories-over-commutative-rings.q001
		\[
			\Hom_{\mathcal{A}}\left(\varinjlim \alpha,T\right)
			\simeq \varprojlim_{i \in \Obj(I)}
			\Hom_{\mathcal{A}}(\alpha(i),T),
			\quad T \in \Obj(\mathcal{A})
		\]
		is $\Bbbk$-linear whenever $\varinjlim \alpha$ exists. The same statement
		holds for $\varprojlim$.
	\end{enumerate}
\end{proposition}

% segment-id: o014.li2.chapter1.linear-categories-over-commutative-rings.p004
Unless otherwise specified, all functors and equivalences between
$\Bbbk\dcate{Mod}$-categories will henceforth be understood to be
$\Bbbk$-linear.

% segment-id: o014.li2.chapter1.linear-categories-over-commutative-rings.d002
\begin{definition}\label{def:k-linear-cat}
	\index{category!$\Bbbk$-linear ($\Bbbk$-linear)}
	If an additive category $\mathcal{A}$ is equipped with a $\Bbbk$-linear
	structure compatible with its existing $\cate{Ab}$-category structure, then
	$\mathcal{A}$ together with these data is called a
	\emph{$\Bbbk$-linear category}.
\end{definition}

% segment-id: o014.li2.chapter1.linear-categories-over-commutative-rings.p005
This amounts to replacing the $\cate{Ab}$-category by a
$\Bbbk\dcate{Mod}$-category in the discussion of \S\ref{sec:Ker-Coker}. Most
results about additive categories in this chapter extend to the $\Bbbk$-linear
setting by exactly the same arguments. The sole exception concerns
Corollary~\ref{prop:automatic-additivity}: in an additive category, addition on
the $\Hom$ sets can be characterized by properties of the category itself, but
this is not true of scalar multiplication by $\Bbbk$. Consequently,
functors---and even equivalences---between $\Bbbk$-linear categories are not
automatically $\Bbbk$-linear.

% segment-id: o014.li2.chapter1.linear-categories-over-commutative-rings.p006
We next collect a few results concerning $\Bbbk$-linearity. Recall that the
\emph{center} of a category $\mathcal{A}$ is defined by
\index{center} \index[sym1]{ZA@$Z(\mathcal{A})$}
$Z(\mathcal{A}) := \End(\identity_{\mathcal{A}})$. Its elements are
endomorphisms of the identity functor $\identity_{\mathcal{A}}$, written as
$(a_X)_{X \in \Obj(\mathcal{A})}$ with
$a_X \in \End_{\mathcal{A}}(X)$. Composition makes the center a commutative
monoid; see \cite[Definition 2.3.8,
Proposition 2.3.9]{Li1}. If $\mathcal{A}$ is an $\cate{Ab}$-category, then
$Z(\mathcal{A})$ becomes a commutative ring under composition and addition
$(a_X)_X + (b_X)_X := (a_X+b_X)_X$.

% segment-id: o014.li2.chapter1.linear-categories-over-commutative-rings.pr002
\begin{proposition}\label{prop:k-linear-center}
	Let $\mathcal{A}$ be an $\cate{Ab}$-category. Then there is a bijection
	% segment-id: o014.li2.chapter1.linear-categories-over-commutative-rings.q002
	\[
		\left\{\text{$\Bbbk\dcate{Mod}$-category structures on $\mathcal{A}$}\right\}
		\stackrel{1:1}{\longleftrightarrow}
		\Hom_{\text{rings}}(\Bbbk,Z(\mathcal{A})).
	\]
	The map is given as follows. If a $\Bbbk\dcate{Mod}$-category structure on
	$\mathcal{A}$ has been fixed, define, for every $a \in \Bbbk$ and
	$X \in \Obj(\mathcal{A})$,
	$a_X := a \cdot \identity_X \in \End_{\mathcal{A}}(X)$. Then
	$a \mapsto (a_X)_{X \in \Obj(\mathcal{A})} \in Z(\mathcal{A})$ is a ring
	homomorphism. Conversely, given a ring homomorphism on the right,
	$a \mapsto (a_X)_{X \in \Obj(\mathcal{A})}$, define a
	$\Bbbk\dcate{Mod}$-category structure on $\mathcal{A}$ by
	$a \cdot f := f \circ a_X = a_Y \circ f$, for $f \in \Hom(X,Y)$.
\end{proposition}

% segment-id: o014.li2.chapter1.linear-categories-over-commutative-rings.pf001
\begin{proof}
	The condition for $(a_X)_{X \in \Obj(\mathcal{A})}$ to belong to
	$Z(\mathcal{A})$ is precisely
	$f \circ a_X = a_Y \circ f$ for every $f \in \Hom(X,Y)$. If
	$\mathcal{A}$ is equipped with a $\Bbbk\dcate{Mod}$-category structure,
	define $a_X := a \cdot \identity_X$. Bilinearity gives
	\[
		f \circ a_X = f \circ (a \cdot \identity_X)
		= (a f) \circ \identity_X
		= (a \cdot \identity_Y) \circ f
		= a_Y \circ f,
	\]
	so $(a_X)_{X \in \Obj(\mathcal{A})} \in Z(\mathcal{A})$.
	The remaining verifications are immediate.
\end{proof}

% segment-id: o014.li2.chapter1.linear-categories-over-commutative-rings.p007
Consequently, the $\cate{Ab}$-category $\mathcal{A}$ is naturally a
$Z(\mathcal{A})\dcate{Mod}$-category.
